\section{Signed coefficient-cone preservation by the actual owner division}

The tangent-scaling chart needed for a relative-to-point comparison has a
SIGNED coefficient weight: the exceptional exponent minus the sum of selected
tangent exponents. Natural nonnegative coordinate-weight bounds do not cover
this cone. The following actual coefficient induction uses the existing
whole monic vertex family and genuine owner recursion, with no division
solution or output cone certificate supplied.

\begin{theorem}[The actual whole family division preserves genuine signed additive coefficient bounds]
\label{hir:actual-family-division-signed-coefficient-bound}
\uses{hir:actual-reduced-basis-entire-filtered-division}
\lean{Hironaka.HironakaActualSignedConeDivision.familyDivisionCoefficient_eq_zero_of_signed_weight_lt,Hironaka.HironakaActualSignedConeDivision.familyFormalQuotient_signed_weight_bound,Hironaka.HironakaActualSignedConeDivision.familyFormalRemainder_signed_weight_bound}
\leanok
Let $\nu:\mathbb N^n\to\mathbb Z$ be an actual additive homomorphism and
$c\in\mathbb Z$. Suppose each actual prescribed monic vertex series $B_a$
has zero coefficient at every $\beta$ with $\nu(\beta)<\nu(a)$, and the
actual input $F$ has zero coefficients below $c$. Then the genuine
coefficient-recursive division coefficients vanish below $c$, every literal
owner quotient has zero coefficient at $\beta$ if $\nu(\beta)+\nu(a)<c$,
and the actual entire remainder vanishes below $c$. Signed weights are
allowed to have negative coordinate values.
\end{theorem}
\begin{proof}
Induct in the genuine well-founded monomial order used by the actual
coefficient recursion, not in the possibly negative signed weight.
At an exponent $\alpha$ below $c$, the original input coefficient is zero.
A nonzero recursive tail summand comes from a genuine antidiagonal pair
$\lambda+\beta=\alpha$ and a vertex $a$. The actual vertex support
hypothesis gives $\nu(a)\le\nu(\beta)$. Additivity therefore gives
$\nu(\lambda+a)\le\nu(\alpha)<c$. This recursive argument is strictly
smaller in the genuine monomial order by the literal tail definition, so
its coefficient is zero by induction. Every finite tail summand is zero,
and the recursion gives the asserted coefficient zero.

The genuine owner quotient at $\beta$ uses precisely the recursive
coefficient at $\beta+a$ when the actual owner is $a$, and otherwise is
zero. Additivity and the preceding induction give the shifted bound.
The actual remainder uses the same recursive coefficient off the entire
initial diagram and zero on it, giving its bound. No positivity or
well-foundedness of $\nu$ has been invoked.
\end{proof}

\begin{theorem}[Every actual zero-direction blowup image has its genuine signed support cone]
\label{hir:actual-zero-chart-signed-support-cone}
\uses{hir:actual-formal-monomial-blowup-exponent-embedding}
\lean{Hironaka.HironakaActualSignedConeDivision.actualBlowupConeWeight,Hironaka.HironakaActualSignedConeDivision.actualFormalBlowup_zero_coefficient_outside_cone}
\leanok
For a prescribed actual chart with pivot $e$ and selected coordinates $A$,
put $\nu(\alpha)=\alpha_e-\sum_{i\in A\setminus\{e\}}\alpha_i$.
Every genuine whole zero-direction blowup image has zero coefficient
where $\nu(\alpha)<0$. This is an output from the actual substitution,
without an assumed image or support formula.
\end{theorem}
\begin{proof}
The literal actual monomial exponent map adds the sum of selected nonpivot
exponents to the exceptional exponent and leaves every other exponent
unchanged. The proved true coordinate-product formula gives this actual
monomial image. Its signed weight is exactly the original nonnegative
exceptional exponent. In the genuine whole substitution coefficient
formula, every summand at a negative-weight exponent is therefore zero:
its coordinate-product monomial cannot have that exponent. The actual
coefficient finitary sum vanishes, proving the whole output support bound.
\end{proof}

\begin{theorem}[Actual image-family owner quotients and remainders retain the signed chart cone]
\label{hir:actual-image-family-owner-division-signed-cone}
\uses{hir:actual-family-division-signed-coefficient-bound,hir:actual-zero-chart-signed-support-cone}
\lean{Hironaka.HironakaActualSignedConeDivision.actual_image_family_division_cone}
\leanok
Suppose an existing actual whole monic vertex family consists of specified
actual zero-chart pullback series, and every one of its actual vertex labels
has signed weight zero. For ANY original formal input series, dividing its
actual zero-chart image by this same family gives literal owner quotients
and an actual remainder supported in the nonnegative chart cone.
The actual image-family equations and zero-weight labels are explicit natural
input data; the quotient and remainder support bounds are outputs.
\end{theorem}
\begin{proof}
The preceding unconditional actual-substitution theorem gives the cone
support of every specified family member and of the actual input image.
Since every vertex has signed weight zero, these are exactly the natural
vertex and input support hypotheses of the first theorem at $c=0$.
Its genuine recursive quotient and remainder bounds therefore prove the
claim for this SAME family and input. No output quotient support condition
or actual division solution is assumed.
\end{proof}

This establishes a genuine signed-cone operator estimate. Producing the
required image monic family from an original relative retained-order chart,
proving the converse actual range-cone dictionary and its analytic inverse
convergence, and completing the relative maximal-filtered syzygy/HS criterion
are separate obligations, not conclusions of this support estimate.
